\chapter{기술 부록 읽는 법}
\begin{warnbox}[먼저 읽기]
이 부록은 따로 만들었던 기술 문서 네 개를 합친 것이다. 기법 정리(피처, 데이터와 검증 기법, 수익률, 모델, 평가 지표), 코드 해설(핵심 코드를 한 줄씩), 개념 치트시트, 코드 치트시트다. 작성 시점이 다르므로(2026년 11월 9--16일) \textbf{최종 제안과 test 판정은 본문(핵심 정리, 9장, 10장)이 우선한다.} 부록에 나오는 1-SE 규칙과 ``부도확률 순 80\%''는 이전 최종안(규칙서 v2.2)의 설명이다. 현재 최종 제안은 기대수익 순(예산 무관), 주 판정은 v1의 H2다(규칙서 v2.3).
\end{warnbox}
\begin{table}[H]\centering\small
\begin{tabularx}{\linewidth}{@{}p{45mm}L@{}}\toprule
부분 & 내용\\\midrule
기법 정리 & 피처 101개의 뜻과 처리, 시점 학습과 결측, 수익률과 비교 기준, 모델과 보정, 평가 지표(AUC, Brier, Sharpe 등), 검증 방법론\\
코드 해설 & \ft{src/lc/}와 \ft{scripts/}의 핵심 코드를 한 줄씩: 데이터, 수익률(IRR), 모델, 선별 곡선, 포트폴리오 Sharpe와 부트스트랩\\
개념 치트시트 & 신용 대출 선별 전략을 검증하는 개념을 두 단 한 쪽으로\\
코드 치트시트 & 복사해서 쓰는 코드 조각(\ft{docs/notes/code/snippets.py})\\\bottomrule
\end{tabularx}
\end{table}
